\subsection{The Last Days and the Vigil}\label{sec:last-days}

From 17 December the season quickens, and the Office marks the days to the Vigil in two ways. At
Lauds of the six ferias before the Vigil the antiphons are proper to each day. At Vespers the
Magnificat takes one of the seven great antiphons, said entire before and after the canticle, one
for each evening to the day before the Vigil. Each calls on Christ by a title and asks him to come,
\latin{veni} (gloss: come):

\begin{center}
\begin{tabular}{@{}>{\raggedright\arraybackslash}p{0.18\linewidth}>{\raggedright\arraybackslash}p{0.34\linewidth}>{\raggedright\arraybackslash}p{0.34\linewidth}@{}}
\toprule
Date & Latin & Bute's English\\
\midrule
17 December & \latin{O Sapientia} & O Wisdom\\
18 December & \latin{O Adonai} & O Adonai\\
19 December & \latin{O radix Jesse} & O Root of Jesse\\
20 December & \latin{O clavis David} & O Key of David\\
21 December & \latin{O Oriens} & O Day-Spring\\
22 December & \latin{O Rex Gentium} & O King of the Gentiles\\
23 December & \latin{O Emmanuel} & O Emmanuel\\
\bottomrule
\end{tabular}
\end{center}

\noindent The full texts, in the Latin of the Breviary and Bute's English, are in
Appendix~\ref{sec:prayers}. In seven evenings the antiphons ask him to teach, to redeem, to deliver, to bring out
the prisoner, to give light, and to save.

\subsubsection{The Vigil}

The Vigil of the Nativity, 24 December, is of the first class. When it falls on the Fourth Sunday
of Advent, it takes that Sunday's place, and the Sunday is not commemorated (n.~30). Advent ends
at None of the Vigil (n.~71). The station is at St Mary Major. The Mass opens \latin{Hodie
scietis} (gloss: this day you shall know), from Exodus 16:6--7, and reads the opening of St Paul's
letter to the Romans (1:1--6) and the Gospel of the angel's message to Joseph (Matthew 1:18--21).
Its Collect:

\begin{quote}
\latin{Deus, qui nos redemptionis nostrae annua exspectatione laetificas: praesta; ut Unigenitum
tuum, quem Redemptorem laeti suscipimus, venientem quoque judicem securi videamus, Dominum
nostrum Jesum Christum, Filium tuum: Qui tecum vivit et regnat in unitate.}\par
{\smallskip\noindent\small ---Collect of the Vigil of the Nativity. Latin: Triptych's editorial
text of the 1962 proper, the wording of the Pustet Missal of 1862 read against the 1962 typical
edition.\par}
\end{quote}

\noindent God gladdens us \latin{annua exspectatione} (gloss: by the yearly expectation) of our
redemption, and the prayer asks that the Son we receive with joy as Redeemer we may also see
\latin{venientem quoque judicem securi} (gloss: without fear, when he comes as judge).

At Matins of the Vigil the Breviary reads St Jerome on the Gospel, from the first book of his
commentary on St Matthew. Of Joseph's resolve to put Mary away privately, he offers this reading:

\begin{quote}
Or this may be considered a testimony to Mary, that Joseph, confident in her purity, and
wondering at what had happened, covered in silence that mystery which he could not explain.\par
{\smallskip\noindent\small ---Jerome, \textit{Commentary on Matthew}, book I, on Matthew 1:19,
as the \textit{Catena Aurea} gives it, Oxford tr. (1842).\par}
\end{quote}

\noindent The angel then tells Joseph what the Child's name will mean:

\begin{quote}
And she shall bring forth a son: and thou shalt call his name Jesus. For he shall save his people
from their sins.\par
{\smallskip\noindent\small ---Matthew 1:21, Douay-Rheims (Challoner).\par}
\end{quote}

Under the Code of Canon Law of 1917 the Vigil was a day of fast and abstinence for the Latin
Church (can.~1252 \S\,2); in 2026 neither the Code of 1983 nor, in the United States, the bishops'
norms of 1966 binds anyone to fast on it. Section~\ref{sec:fasting} gives the law in full.

\medskip
\noindent\textit{Counsel for the last days.} Say each evening's O antiphon at the end of the day's
prayer, from 17 to 23 December (Appendix~\ref{sec:prayers}). On the Vigil read Matthew 1:18--21
and Jerome's sentence on Joseph, and keep some part of the day in silence before the mystery, as
Joseph did. Those who wish to keep the Vigil fast of the older law may do so as a devotion, within
what health and duties allow.
